\section{Klassen \& OOP}\label{sec:02-klassen}

\subsection{Klasse definieren}\label{sec:02-klasse-definieren}
\ZSFindex[oop]{OOP / Objektorientierte Programmierung}
\begin{codebox}[Klasse definieren und verwenden]
\begin{lstlisting}[style=CodeExpert]
class Car:
    def __init__(self, brand, model):
        self.brand = brand
        self.model = model

    def honk(self):
        return f"{self.model} is honking!"

    def announce(self):
        return self.honk()  # eigene Methode aufrufen

car = Car("Volvo", "EX30")  # car wird in __init__ zu self
message = car.honk()  # car wird automatisch als self übergeben
model = car.model  # Attribut des Objekts lesen
\end{lstlisting}
\end{codebox}

\begin{defbox}[\texttt{self} und Methodenaufruf]
\begin{factlist}
\ZSFFact{Methode definieren}{\texttt{def method(self, argument)} — \texttt{self} bezeichnet das aktuelle Objekt.}
\ZSFFact{Von aussen aufrufen}{\texttt{object.method(argument)} — das Objekt wird automatisch als \texttt{self} übergeben.}
\ZSFFact{Innerhalb aufrufen}{\texttt{self.method(argument)} — Methode desselben Objekts verwenden.}
\ZSFFact{Zustand speichern}{\texttt{self.attribute} bleibt im Objekt gespeichert; eine lokale Variable nur während des Aufrufs.}
\end{factlist}
\end{defbox}

\subsection{Vererbung}\label{sec:02-vererbung}
\ZSFindex{Vererbung}
\begin{defbox}[Oberklasse / Unterklasse]
Eine \ZSFkeyword*{Oberklasse} (Superclass) ist eine Klasse, die als Basis für eine andere Klasse dient – die sogenannte \ZSFkeyword*{Unterklasse} (Subclass). Mittels \verb|super().__init__()| kann die Unterklasse auf den Konstruktor der Oberklasse zugreifen.
\end{defbox}
\begin{codebox}[super() in der Unterklasse]
\begin{lstlisting}[style=CodeExpert]
class ElectricCar(Car):  # Erbt alle Attribute & Methoden von Car
    def __init__(self, brand, model):
        super().__init__(brand, model)  # Car.__init__ erzeugt self.brand & self.model

    def charge(self):
        return f"{self.model} battery is charging."  # Greift auf self.model aus Car zu
\end{lstlisting}
\end{codebox}

\subsection{Dunder-Methoden \& Operator-Überladung}\label{sec:02-klasseninterne-funktionen-implementieren}
\ZSFindex{Dunder-Methode}
\ZSFindex[operatoruberladung]{Operatorüberladung}
\begin{statementbox}[Rückgabewert]
\texttt{\_\_add\_\_} bestimmt das Ergebnis von \texttt{objekt\_a + objekt\_b}. Der Rückgabewert muss zur Bedeutung des Operators passen: Häufig ist er eine neue Instanz derselben Klasse, andere Ergebnistypen sind aber ebenfalls möglich.
\end{statementbox}
\begin{codebox}[\texttt{\_\_add\_\_} (Beispiel: Reissverschluss-Logik)]
\begin{lstlisting}[style=CodeExpert]
class CustomString:
    def __init__(self, string):
        self.string = string

    def __str__(self):  # Bestimmt Textausgabe bei print(objekt)
        return self.string

    def __add__(self, other):  # self = links (s1), other = rechts (s2)
        result_string = ""
        for i in range(len(self.string)):
            if i < len(other.string):
                result_string += (self.string[i]
                                  + other.string[i])
            else:
                result_string += self.string[i]
        return CustomString(result_string)  # Gibt neues Objekt zurück

s1 = CustomString("abcd")
s2 = CustomString("123")  # s1 + s2 ruft s1.__add__(s2) auf
print(s1 + s2) # -> a1b2c3d
\end{lstlisting}
\end{codebox}



%=====================================================================================================================

%=====================================================================================================================

%=====================================================================================================================

%=====================================================================================================================

    


%=====================================================================================================================


%%%%%%%%%%%%%%%%%%%%%%%%%%%%%%%%%%%%%%%%%%%%%%%%%%%%%%%%%%%%%%%%%%%%%%%%%%%%%%%%%%%%%%%%%%%%%%%%%%%%%%%%%%%%%%%%%%%%%%
